\answer{ex:14:1}{(a)~$0.17$~s와 $1.5$~s. (b)~$11$~cm보다 짧은 거리부터.}
\answer{ex:14:2}{$g\le\beta w_{q\mu}/w_{z\mu}=5$인 동안 통증 없는 촉각은 어떤 $\mu$에서도 통과하지 못한다; 촉각 $\mu=1$은 $g>6$에서 통과한다: 강한 감작은 보통의 접촉을 통증으로 만든다.}
\answer{ex:14:3}{(a)~$g^*=(1-k_sC)/(1-k_sA)$, $k_sA<1$이면 안정. (b)~촉각 없이는 $\nu\ge2$에서, 촉각이 있으면 이미 $\nu\ge1.7$에서.}
\answer{ex:14:4}{$V(t)=-Pe^{-(T-t)/\tau_\gamma}$. (a)~$-Pe^{-T/\tau_\gamma}\approx-0.37P$. (b)~$+Pe^{-(T-t_e)/\tau_\gamma}\approx+0.61P$; $t_e$와 함께 $+P$까지 커지며, 통증을 피하도록 가르친다.}
\answer{ex:14:5}{$k_s=4$에서는 걸림이 없다(둘째 안정 상태는 $k_s\ge4.58$에서 존재); $k_s=4.6$, $5$, $6$, $8$에서 $T_{\min}\approx4.35$, $1.40$, $0.65$, $0.32\,\tau_s$.}
\answer{ex:14:6}{$w_{q\nu}=4/9\approx0.444$, $q_0=2/9\approx0.222$, $w_{q\mu}=49/45\approx1.089$; 통증 없는 촉각은 $g=49/9\approx5.4$까지 안전하다.}
\answer{ex:14:7}{(a)~문턱값은 $\nu_c\ge q_0/w_{q\nu}$이면 $\nu_c(g)=\theta/g$, 아니면 $(\theta+\beta q_0)/(g+\beta w_{q\nu})$; $g^*\approx2.45$, $1.0$, $0.25$. (b)~열린 문제.}
